% GENERATED FILE. Edit canonical lessons, then run npm run generate:books.
% canonical-source-hash: fnv1a64:6482babd5d516fe1
% canonical-lessons: AR-C347-ikhtatafa, AR-C347-lahiqa, AR-C347-rahina, AR-C347-fidya, AR-C347-khatif

\chapter{A kidnapping}
\label{ch:ar-a-kidnapping}
\hlchaptermodality{347}

\begin{chapteropening}
\textbf{By the end of this chapter:} \emph{I can say to kidnap, to catch up with, a hostage, a ransom and a kidnapper.}

Complete the last lesson of chapter 347: I can say to kidnap, to catch up with, a hostage, a ransom and a kidnapper.
\end{chapteropening}

\section[\texorpdfstring{\ar{اختطف} (ikhtaṭafa)}{اختطف (ikhtaṭafa)}]{\ar{اختطف} (ikhtaṭafa) — to kidnap; to snatch}

\label{lesson:AR-C347-ikhtatafa}



\begin{quote}
Before the new one: say the Arabic for to sneak in, then the Arabic for to chase.
\end{quote}

\subsection*{You'll want to know: \ar{اختطف}}
\textbf{\ar{اختطف}} — \emph{ikhtaṭafa} — \textquotedblleft{}to kidnap; to snatch\textquotedblright{}.

\begin{quote}
\ar{اختطف} \ar{رجلان} \ar{ابن} \ar{التاجر}.
\end{quote}

\begin{quote}
\emph{ikhtaṭafa rajulān ibn at-tājir.}
\end{quote}

\begin{quote}
Two men kidnapped the merchant’s son.
\end{quote}

Present: \ar{يختطف}, \emph{yakhtaṭif}.

\begin{cousinweb}
\textbf{\ar{اختطف}} (\emph{ikhtaṭafa}) is built on the native Arabic root \emph{kh-ṭ-f} (Lane), \textquotedblleft{}to snatch away quickly\textquotedblright{}, as a bird snatches its prey. It is Form VIII.
\end{cousinweb}

\begin{grammarlens}[title={What you've built}]
One more word for telling what happened.
\end{grammarlens}

\subsection*{Your turn}
\begin{itemize}
\raggedright
  \item \emph{Say it:} \emph{ikhtaṭafa}
  \item \emph{Say it:} \emph{ikhtaṭafa}, once more
  \item \emph{Say it:} \emph{ikhtaṭafa rajulān ibn at-tājir.} — Two men kidnapped the merchant’s son.
  \item \emph{Recall:} say the Arabic for to sneak in, then the Arabic for to chase, then say \emph{ikhtaṭafa} again
\end{itemize}

\begin{tcolorbox}[breakable,colback=teal!4,colframe=teal!35!black,title={Before you move on}]
What does \ar{اختطف} mean? (\textquotedblleft{}To kidnap; to snatch\textquotedblright{}.) Say it once more.
\end{tcolorbox}
\section[\texorpdfstring{\ar{لحق} (laḥiqa)}{لحق (laḥiqa)}]{\ar{لحق} (laḥiqa) — to catch up with}

\label{lesson:AR-C347-lahiqa}



\begin{quote}
Before the new one: say the Arabic for to chase, then the Arabic for to kidnap.
\end{quote}

\subsection*{You'll want to know: \ar{لحق}}
\textbf{\ar{لحق}} — \emph{laḥiqa} — \textquotedblleft{}to catch up with\textquotedblright{}.

\begin{quote}
\ar{لحق} \ar{الشرطي} \ar{باللص} \ar{في} \ar{الشارع}.
\end{quote}

\begin{quote}
\emph{laḥiqa ash-shurṭī bi-l-liṣṣ fī ash-shāriʿ.}
\end{quote}

\begin{quote}
The policeman caught up with the thief in the street.
\end{quote}

Present: \ar{يلحق}, \emph{yalḥaq}.

\ar{لحق} \ar{بالقطار}: he caught the train.

\begin{cousinweb}
\textbf{\ar{لحق}} (\emph{laḥiqa}) is built on the native Arabic root \emph{l-ḥ-q} (Lane), \textquotedblleft{}to overtake, to come up with\textquotedblright{}, the root of the taught \emph{lāḥiqan}, later.
\end{cousinweb}

\begin{grammarlens}[title={What you've built}]
One more word for telling what happened.
\end{grammarlens}

\subsection*{Your turn}
\begin{itemize}
\raggedright
  \item \emph{Say it:} \emph{laḥiqa}
  \item \emph{Say it:} \emph{laḥiqa}, once more
  \item \emph{Say it:} \emph{laḥiqa ash-shurṭī bi-l-liṣṣ fī ash-shāriʿ.} — The policeman caught up with the thief in the street.
  \item \emph{Recall:} say the Arabic for to chase, then the Arabic for to kidnap, then say \emph{laḥiqa} again
\end{itemize}

\begin{tcolorbox}[breakable,colback=teal!4,colframe=teal!35!black,title={Before you move on}]
What does \ar{لحق} mean? (\textquotedblleft{}To catch up with\textquotedblright{}.) Say it once more.
\end{tcolorbox}
\section[\texorpdfstring{\ar{رهينة} (rahīna)}{رهينة (rahīna)}]{\ar{رهينة} (rahīna) — a hostage}

\label{lesson:AR-C347-rahina}



\begin{quote}
Before the new one: say the Arabic for to kidnap, then the Arabic for to catch up with.
\end{quote}

\subsection*{You'll want to know: \ar{رهينة}}
\textbf{\ar{رهينة}} — \emph{rahīna} — \textquotedblleft{}a hostage\textquotedblright{}.

\begin{quote}
\ar{بقيت} \ar{الرهينة} \ar{ثلاثة} \ar{أيام} \ar{في} \ar{البيت}.
\end{quote}

\begin{quote}
\emph{baqiyat ar-rahīna thalāthat ayyām fī al-bayt.}
\end{quote}

\begin{quote}
The hostage stayed in the house for three days.
\end{quote}

\begin{cousinweb}
\textbf{\ar{رهينة}} (\emph{rahīna}) is built on the native Arabic root \emph{r-h-n} (Lane), \textquotedblleft{}to give a thing as a pledge\textquotedblright{}; \emph{rahn} is a pawn. A \emph{rahīna} is a person held as one.
\end{cousinweb}

\begin{grammarlens}[title={What you've built}]
One more word for telling what happened.
\end{grammarlens}

\subsection*{Your turn}
\begin{itemize}
\raggedright
  \item \emph{Say it:} \emph{rahīna}
  \item \emph{Say it:} \emph{rahīna}, once more
  \item \emph{Say it:} \emph{baqiyat ar-rahīna thalāthat ayyām fī al-bayt.} — The hostage stayed in the house for three days.
  \item \emph{Recall:} say the Arabic for to kidnap, then the Arabic for to catch up with, then say \emph{rahīna} again
\end{itemize}

\begin{tcolorbox}[breakable,colback=teal!4,colframe=teal!35!black,title={Before you move on}]
What does \ar{رهينة} mean? (\textquotedblleft{}A hostage\textquotedblright{}.) Say it once more.
\end{tcolorbox}
\section[\texorpdfstring{\ar{فدية} (fidya)}{فدية (fidya)}]{\ar{فدية} (fidya) — a ransom}

\label{lesson:AR-C347-fidya}



\begin{quote}
Before the new one: say the Arabic for to catch up with, then the Arabic for a hostage.
\end{quote}

\subsection*{You'll want to know: \ar{فدية}}
\textbf{\ar{فدية}} — \emph{fidya} — \textquotedblleft{}a ransom\textquotedblright{}.

\begin{quote}
\ar{طلبت} \ar{العصابة} \ar{فدية} \ar{كبيرة}.
\end{quote}

\begin{quote}
\emph{ṭalabat al-ʿiṣāba fidya kabīra.}
\end{quote}

\begin{quote}
The gang demanded a large ransom.
\end{quote}

\begin{cousinweb}
\textbf{\ar{فدية}} (\emph{fidya}) is built on the native Arabic root \emph{f-d-y} (Lane), \textquotedblleft{}to ransom, to give a thing for another\textquotedblright{}; \emph{fidāʾ} is the same.
\end{cousinweb}

\begin{grammarlens}[title={What you've built}]
One more word for telling what happened.
\end{grammarlens}

\subsection*{Your turn}
\begin{itemize}
\raggedright
  \item \emph{Say it:} \emph{fidya}
  \item \emph{Say it:} \emph{fidya}, once more
  \item \emph{Say it:} \emph{ṭalabat al-ʿiṣāba fidya kabīra.} — The gang demanded a large ransom.
  \item \emph{Recall:} say the Arabic for to catch up with, then the Arabic for a hostage, then say \emph{fidya} again
\end{itemize}

\begin{tcolorbox}[breakable,colback=teal!4,colframe=teal!35!black,title={Before you move on}]
What does \ar{فدية} mean? (\textquotedblleft{}A ransom\textquotedblright{}.) Say it once more.
\end{tcolorbox}
\section[\texorpdfstring{\ar{خاطف} (khāṭif)}{خاطف (khāṭif)}]{\ar{خاطف} (khāṭif) — a kidnapper}

\label{lesson:AR-C347-khatif}



\begin{quote}
Before the new one: say the Arabic for a hostage, then the Arabic for a ransom.
\end{quote}

\subsection*{You'll want to know: \ar{خاطف}}
\textbf{\ar{خاطف}} — \emph{khāṭif} — \textquotedblleft{}a kidnapper\textquotedblright{}.

\begin{quote}
\ar{هرب} \ar{الخاطف} \ar{بالسيارة}.
\end{quote}

\begin{quote}
\emph{haraba al-khāṭif bi-s-sayyāra.}
\end{quote}

\begin{quote}
The kidnapper got away in the car.
\end{quote}

Plural \ar{خاطفون}. As an adjective it is swift: \ar{زيارة} \ar{خاطفة}, a flying visit.

\begin{cousinweb}
\textbf{\ar{خاطف}} (\emph{khāṭif}) is built on the native Arabic root \emph{kh-ṭ-f} (Lane), \textquotedblleft{}to snatch\textquotedblright{}. \emph{Khāṭif} is the active participle, \textquotedblleft{}one who snatches\textquotedblright{}.
\end{cousinweb}

\begin{grammarlens}[title={What you've built}]
One more word for telling what happened.
\end{grammarlens}

\subsection*{Your turn}
\begin{itemize}
\raggedright
  \item \emph{Say it:} \emph{khāṭif}
  \item \emph{Say it:} \emph{khāṭif}, once more
  \item \emph{Say it:} \emph{haraba al-khāṭif bi-s-sayyāra.} — The kidnapper got away in the car.
  \item \emph{Recall:} say the Arabic for a hostage, then the Arabic for a ransom, then say \emph{khāṭif} again
\end{itemize}

\begin{tcolorbox}[breakable,colback=teal!4,colframe=teal!35!black,title={Before you move on}]
What does \ar{خاطف} mean? (\textquotedblleft{}A kidnapper\textquotedblright{}.) Say it once more.
\end{tcolorbox}
